% ---------------------------------------------------------------------------- %
%% \file rdmpafsym4a.tex
%% \author Sebastien Beaugrand
%% \sa http://beaugrand.chez.com/
%% \copyright CeCILL 2.1 Free Software license
%% \brief Poutre sur quatre appuis avec porte-à-faux symétrique
% ---------------------------------------------------------------------------- %
\documentclass[a4paper]{article}
\usepackage{vmargin}
\setmarginsrb{10mm}{10mm}{10mm}{10mm}{0mm}{0mm}{0mm}{0mm}
\usepackage[french]{babel}
\usepackage[utf8]{inputenc}
\usepackage[T1]{fontenc}
\usepackage{EPB_SI/EPB_SI}
\pagestyle{empty}
\parindent0cm

\begin{document}

% ---------------------------------------------------------------------------- %
% Effort tranchant et moment fléchissant
% ---------------------------------------------------------------------------- %
\begin{minipage}{9.5cm}
\centerline{\textbf{Poutre chargée uniformément sur}}
\centerline{\textbf{quatre appuis avec porte-à-faux symétrique}}
~\\
~\\
Appuis:
$$R_A = R_D = p\left(a + \frac{b}{2}\right) + \frac{M_B - M_A}{b}$$
$$R_B = R_C = \frac{p(b+l)}{2} - \frac{M_B - M_A}{b}$$
~\\
Vérification: $2R_A + 2R_B = p(2a + 2b + l)$\\
~\\
~\\
\small
\begin{tabular}{|c|c|c|}\hline
$x$&
Effort tranchant&
Moment fléchissant\\
\hline
$0$&
$0$&
$0$\\
\hline
$a^-$&
$pa$&
\\
\hline
$a^+$&
$pa - R_A$&
$-p\frac{a^2}{2}$\\
\hline
$a+b^-$&
$pa - R_A + pb$&
$$\\
\hline
$a+b^+$&
$\ldots - R_B = -p\frac{l}{2}$&
$-p\frac{a^2}{2}-pab+R_Ab-p\frac{b^2}{2}$\\
\hline
$a+b+\frac{l}{2}$&
$-p\frac{l}{2} + p\frac{l}{2} = 0$&
$M_B + p\frac{l}{2}\frac{l}{4} = M_B + p\frac{l^2}{8}$\\
\hline
\end{tabular}
\normalsize
\end{minipage}
\begin{minipage}{9.5cm}
\begin{tikzpicture}
\def\L{6.0}
\def\a{1.0}
\def\b{2.0}
\def\l{2.0}
\def\c{4.0}
\def\d{5.0}
% ---------------------------------------------------------------------------- %
% Tracé de la poutre
% ---------------------------------------------------------------------------- %
\PoutreAppuiSimple{\a}{0}
\PoutreAppuiSimple{\b}{0}
\PoutreAppuiSimple{\c}{0}
\PoutreAppuiSimple{\d}{0}
\PoutreBaseLocale{\L}{0.5}
\PoutreChargeRepartie{0}{0}{\L}{$p$}
\node at (\a+0.2,0)[below right] {$A$};
\node at (\b+0.2,0)[below right] {$B$};
\node at (\c+0.2,0)[below right] {$C$};
\node at (\d+0.2,0)[below right] {$D$};
\draw[line width=2.5pt] (0,0) -- (\L,0);
\node at (0,0)[above, left=0.2] {$O$};
\draw ( 0,-0.1) -- ( 0,-1.3);
\draw (\a,-0.1) -- (\a,-1.3);
\draw (\b,-0.1) -- (\b,-1.3);
\draw (\c,-0.1) -- (\c,-1.3);
\draw (\d,-0.1) -- (\d,-1.3);
\draw (\L,-0.1) -- (\L,-1.3);
\draw[<->,>=latex] ( 0,-1.1) -- (\a,-1.1) node [midway, above] {$a$};
\draw[<->,>=latex] (\a,-1.1) -- (\b,-1.1) node [midway, above] {$b$};
\draw[<->,>=latex] (\b,-1.1) -- (\c,-1.1) node [midway, above] {$l$};
\draw[<->,>=latex] (\c,-1.1) -- (\d,-1.1) node [midway, above] {$b$};
\draw[<->,>=latex] (\d,-1.1) -- (\L,-1.1) node [midway, above] {$a$};
~\\
~\\
~\\
% ---------------------------------------------------------------------------- %
% Diagrammes des efforts intérieurs
% ---------------------------------------------------------------------------- %
\begin{scope}[yshift=-3cm]
\PoutreDiagAxesCfg[red]
% Ma = -1/2
% Mb = -7/32
% Ra = 3/2 - 7/32 + 1/2 = 57/32
% Rb = 3/2 + 7/32 - 1/2 = 39/32
% ---------------------------------------------------------------------------- %
% Diagramme Ty
% ---------------------------------------------------------------------------- %
\def\scaleT{0.6}
\begin{scope}
\filldraw[diagCourbe]
 (0,0) --
 plot [smooth,domain= 0:\a] (\x,{\scaleT*\x}) --
 plot [smooth,domain=\a:\b] (\x,{\scaleT*(\x-\a-57/32+\a)}) --
 plot [smooth,domain=\b:\c] (\x,{\scaleT*(\x-\b-\l/2)}) --
 plot [smooth,domain=\c:\d] (\x,{\scaleT*(\x-\d+57/32-\a)}) --
 plot [smooth,domain=\d:\L] (\x,{\scaleT*(\x-\L)}) --
 cycle;
\PoutreDiagAxes{$T_y$}{\L}{-1}{1}
\draw (-0.05,-\scaleT*\l/2) node[right]{\small{$-p\frac{l}{2}$}} --
      ( 0.05,-\scaleT*\l/2);
\draw (-0.1, \scaleT*\a) node[left]{\small{$pa$}} --
      ( 0.1, \scaleT*\a);
\draw (-0.1, \scaleT*\a-\scaleT*57/32) node[left]{\small{$pa-R_A$}} --
      ( 0.1, \scaleT*\a-\scaleT*57/32);
\end{scope}
% ---------------------------------------------------------------------------- %
% Diagramme Mfz
% ---------------------------------------------------------------------------- %
\def\scaleM{1.6}
\begin{scope}[yshift=-3cm]
\filldraw[diagCourbe]
 plot [smooth,domain= 0:\a] (\x,{-\scaleM*\x^2/2)}) --
 plot [smooth,domain=\a:\b] (\x,{-\scaleM*((\x-57/32)^2/2+\a^2/2-(57/32-\a)^2/2)}) --
 plot [smooth,domain=\b:\c] (\x,{-\scaleM*((\x-\b-\l/2)^2/2+7/32-\l^2/8)}) --
 plot [smooth,domain=\c:\d] (\x,{-\scaleM*((\x-\c-7/32)^2/2+7/32)}) --
 plot [smooth,domain=\d:\L] (\x,{-\scaleM*(\x-\L)^2/2}) --
 cycle;
\PoutreDiagAxes{\Mfz}{\L}{-0.8}{1.2}
\draw (-0.1,-\scaleM*7/32) node[left]{\small{$M_B$}} --
      ( 0.1,-\scaleM*7/32);
\draw (-0.1,-\scaleM*\a^2/2) node[left]{\small{$-p\frac{a^2}{2}$}} --
      ( 0.1,-\scaleM*\a^2/2);
\draw (-0.1,-\scaleM*7/32+\scaleM*\l^2/8) node[left]{\small{$M_B+p\frac{l^2}{8}$}} --
      ( 0.1,-\scaleM*7/32+\scaleM*\l^2/8);
\end{scope}
\end{scope}
\end{tikzpicture}
\end{minipage}\\
~\\
~\\
~\\
% ---------------------------------------------------------------------------- %
% Déformée entre A et B
% ---------------------------------------------------------------------------- %
\begin{minipage}[t]{9.5cm}
Déformée entre $A$ et $B$:\\
~\\
$EIy^{''} =
M_A(1-\frac{x}{b}) + M_B\frac{x}{b} + p\frac{x(b-x)}{2}
$\\
~\\
$\phantom{EIy^{''}} =
- p\frac{x^2}{2}
+ (\frac{M_B}{b} - \frac{M_A}{b} + p\frac{b}{2})x
+ M_A
$\\
~\\
$EIy^{'} =
- p\frac{x^3}{6}
+ (\frac{M_B}{b} - \frac{M_A}{b} + p\frac{b}{2})\frac{x^2}{2}
+ M_Ax
+ c_1
$\\
~\\
$EIy =
- p\frac{x^4}{24}
+ (\frac{M_B}{b} - \frac{M_A}{b} + p\frac{b}{2})\frac{x^3}{6}
+ M_A\frac{x^2}{2}
+ c_1x
+ c_2
$\\
~\\
~$
\begin{cases}
  y_{0(x=0)} = 0 \implies c_2 = 0\\~\\
  y_{A(x=b)} = 0
\end{cases}
$\\
~\\
~\\
$c_1 = p\frac{b^4}{24b}
- (\frac{M_B}{b} - \frac{M_A}{b} + p\frac{b}{2})\frac{b^3}{6b}
- M_A\frac{b^2}{2b}
$\\
~\\
$\phantom{c_1} = -p\frac{b^3}{24} - M_B\frac{b}{6} - M_A\frac{b}{3}$\\
~\\
~\\
~\\
% ---------------------------------------------------------------------------- %
% Déformée entre 0 et A
% ---------------------------------------------------------------------------- %
Déformée entre $0$ et $A$:\\
~\\
$EIy^{''} = -p\frac{x^2}{2}$\\
~\\
$EIy^{'} = -p\frac{x^3}{6} + c_3$\\
~\\
$EIy = -p\frac{x^4}{24} + c_3x + c_4$\\
~\\
~$
\begin{cases}
  EIy^{'}_{A(x=a)} = c_1\\~\\
  y_{A(x=a)} = 0
\end{cases}
\begin{cases}
  c_3 = p\frac{a^3}{6} - p\frac{b^3}{24} - M_B\frac{b}{6} - M_A\frac{b}{3}\\~\\
  c_4 = p\frac{a^4}{24} - c_3a
\end{cases}
$\\
~\\
~\\
$c_4 = p\frac{a^4}{24}
- p\frac{a^4}{6}
+ p\frac{ab^3}{24}
+ M_B\frac{ab}{6}
+ M_A\frac{ab}{3}
$\\
~\\
$\phantom{c_4} =
- p\frac{a^4}{8} + p\frac{ab^3}{24} + M_B\frac{ab}{6} + M_A\frac{ab}{3}$\\
~\\
$\phantom{c_4} =
- pi\frac{a^4}{8} + p\frac{ab^3}{24} + M_B\frac{ab}{6} - p\frac{a^3b}{6}$
\end{minipage}
% ---------------------------------------------------------------------------- %
% Moments en A et B
% ---------------------------------------------------------------------------- %
\begin{minipage}[t]{9cm}
Moment en $A$: $M_A = -p\frac{a^2}{2}$\\
~\\
Moment en $B$ avec Clapeyron sous charge uniforme:\\
~\\
$M_{i-1}L_i+2M_i(L_i+L_{i+1})+M_{i+1}L_{i+1}
 = -p\frac{L_i^3}{4}-p\frac{L_{i+1}^3}{4}$\\
~\\
$M_Ab + 2M_B(b+l) + M_Bl = -p\frac{b^3+l^3}{4}$\\
~\\
$M_B = -p\frac{b^3+l^3-2a^2b}{4(2b+3l)}$\\
~\\
Vérification: avec $a = 0$ et $b = l$, on retrouve $M = -p\frac{l^2}{10}$\\
~\\
~\\
% ---------------------------------------------------------------------------- %
% Flèche
% ---------------------------------------------------------------------------- %
Flèche en $E(x=\frac{l}{2})$:\\
~\\
$EIy = -p\frac{l^4}{16\times24}+p\frac{l^4}{16\times6}
 -p\frac{l^4}{16\times3}-M_B\frac{l^2}{8}$\\
~\\
$\phantom{EIy} = -p\frac{5l^4}{16\times24}-M_B\frac{l^2}{8}$\\
~\\
~\\
Pour des flèches en $E$ et $0$ de même grandeur:\\
~\\
$-p\frac{5l^4}{16\times24}-M_B\frac{l^2}{8} \simeq
 -p\frac{a^4}{8} + p\frac{ab^3}{24} + M_B\frac{ab}{6} - p\frac{a^3b}{6}$\\
~\\
Résoudre:\\
\texttt{\footnotesize
def f(l):\\
\mbox{~~a = (L - l) / 2 - b}\\
\mbox{~~m = -(b**3 + l**3 - 2 * a**2 * b) / (2 * b + 3 * l) / 4}\\
\mbox{~~e = -5 * l**4 / 384 - m * l**2 / 8}\\
\mbox{~~e -= -a**4 / 8 + a * b**3 / 24 + m * a * b / 6 - a**3 * b / 6}\\
\mbox{~~return e}\\
}
~\\
Exemple avec $L=1.07$ et $b=0.23$:\\
\texttt{fsolve(f, 0.40)} retourne $0.36$\\
~\\
% ---------------------------------------------------------------------------- %
% pyBar
% ---------------------------------------------------------------------------- %
Déformée avec pyBar:\\
~\\
\includegraphics[scale=0.5]{defo1.pdf}
\end{minipage}
\end{document}
